% GENERATED FILE. Edit canonical lessons, then run npm run generate:books.
% canonical-source-hash: fnv1a64:8f48ff2d43b2de55
% canonical-lessons: FR-C489-fiancailles, FR-C489-alliance, FR-C489-noces, FR-C489-lune-de-miel, FR-C489-ceremonie, FR-R489-revision-482-489

\chapter{Fiançailles et noces}
\label{ch:fr-fiancailles-et-noces}
\hlchaptermodality{489}

\begin{chapteropening}
\textbf{By the end of this chapter:} \emph{I can talk about engagements and weddings.}

Complete the last lesson of chapter 489: I can talk about engagements and weddings.
\end{chapteropening}

\section[les fiançailles]{les fiançailles — engagement (to marry)}

\label{lesson:FR-C489-fiancailles}



\begin{quote}
Before the new one: say the French for to whisper, then the French for to entrust.
\end{quote}

\subsection*{You'll want to know: les fiançailles}
\textbf{les fiançailles} — \textquotedblleft{}engagement (to marry)\textquotedblright{}.

\begin{quote}
Ils ont fêté leurs fiançailles avec toute la famille.
\end{quote}

\begin{quote}
\emph{They celebrated their engagement with the whole family.}
\end{quote}

Always plural and feminine: \emph{des fiançailles}.

\begin{cousinweb}
\emph{les fiançailles} is built on \emph{fiancer} \textquotedblleft{}to promise in marriage\textquotedblright{}, from old \emph{fiance} \textquotedblleft{}trust, promise\textquotedblright{}, which goes back to Vulgar Latin \emph{fīdāre}, from Latin \emph{fīdere} \textquotedblleft{}to trust\textquotedblright{} (Lewis \& Short).
\end{cousinweb}

\begin{grammarlens}[title={What you've built}]
One more word for telling, explaining and describing at B1.
\end{grammarlens}

\subsection*{Your turn}
\begin{itemize}
\raggedright
  \item \emph{Say it:} \emph{les fiançailles}
  \item \emph{Say it:} \emph{les fiançailles}, once more
  \item \emph{Say it:} a sentence of your own with \emph{les fiançailles}
  \item \emph{Recall:} say the French for to whisper, then the French for to entrust, then say \emph{les fiançailles} again
\end{itemize}

\begin{tcolorbox}[breakable,colback=teal!4,colframe=teal!35!black,title={Before you move on}]
What does les fiançailles mean? (\textquotedblleft{}Engagement (to marry)\textquotedblright{}.) Say it once more.
\end{tcolorbox}
\section[l'alliance]{l'alliance — wedding ring; alliance}

\label{lesson:FR-C489-alliance}



\begin{quote}
Before the new one: say the French for to entrust, then the French for engagement.
\end{quote}

\subsection*{You'll want to know: l'alliance}
\textbf{l'alliance} — \textquotedblleft{}wedding ring; alliance\textquotedblright{}.

\begin{quote}
Il porte son alliance à la main gauche.
\end{quote}

\begin{quote}
\emph{He wears his wedding ring on his left hand.}
\end{quote}

Feminine. \emph{Une bague de fiançailles} — an engagement ring.

\begin{cousinweb}
\emph{l'alliance} is built on \emph{allier} \textquotedblleft{}to unite\textquotedblright{}, inherited from Latin \emph{alligāre} \textquotedblleft{}to bind to\textquotedblright{}, built on \emph{ligāre} \textquotedblleft{}to tie\textquotedblright{} (Lewis \& Short).
\end{cousinweb}

\begin{grammarlens}[title={What you've built}]
One more word for telling, explaining and describing at B1.
\end{grammarlens}

\subsection*{Your turn}
\begin{itemize}
\raggedright
  \item \emph{Say it:} \emph{l'alliance}
  \item \emph{Say it:} \emph{l'alliance}, once more
  \item \emph{Say it:} a sentence of your own with \emph{l'alliance}
  \item \emph{Recall:} say the French for to entrust, then the French for engagement, then say \emph{l'alliance} again
\end{itemize}

\begin{tcolorbox}[breakable,colback=teal!4,colframe=teal!35!black,title={Before you move on}]
What does l'alliance mean? (\textquotedblleft{}Wedding ring; alliance\textquotedblright{}.) Say it once more.
\end{tcolorbox}
\section[les noces]{les noces — wedding (celebration)}

\label{lesson:FR-C489-noces}



\begin{quote}
Before the new one: say the French for engagement, then the French for wedding ring.
\end{quote}

\subsection*{You'll want to know: les noces}
\textbf{les noces} — \textquotedblleft{}wedding (celebration)\textquotedblright{}.

\begin{quote}
Mes grands-parents ont fêté leurs noces d'or.
\end{quote}

\begin{quote}
\emph{My grandparents celebrated their golden wedding anniversary.}
\end{quote}

Usually plural. \emph{Le voyage de noces} — the honeymoon trip.

\begin{cousinweb}
\emph{les noces} is inherited from Vulgar Latin \emph{noptiae}, a form of Latin \emph{nūptiae} \textquotedblleft{}a wedding\textquotedblright{}, from \emph{nūbere} \textquotedblleft{}to marry\textquotedblright{} (Lewis \& Short).
\end{cousinweb}

\begin{grammarlens}[title={What you've built}]
One more word for telling, explaining and describing at B1.
\end{grammarlens}

\subsection*{Your turn}
\begin{itemize}
\raggedright
  \item \emph{Say it:} \emph{les noces}
  \item \emph{Say it:} \emph{les noces}, once more
  \item \emph{Say it:} a sentence of your own with \emph{les noces}
  \item \emph{Recall:} say the French for engagement, then the French for wedding ring, then say \emph{les noces} again
\end{itemize}

\begin{tcolorbox}[breakable,colback=teal!4,colframe=teal!35!black,title={Before you move on}]
What does les noces mean? (\textquotedblleft{}Wedding (celebration)\textquotedblright{}.) Say it once more.
\end{tcolorbox}
\section[la lune de miel]{la lune de miel — honeymoon}

\label{lesson:FR-C489-lune-de-miel}



\begin{quote}
Before the new one: say the French for wedding ring, then the French for wedding.
\end{quote}

\subsection*{You'll want to know: la lune de miel}
\textbf{la lune de miel} — \textquotedblleft{}honeymoon\textquotedblright{}.

\begin{quote}
Pour leur lune de miel, ils sont partis au bord de la mer.
\end{quote}

\begin{quote}
\emph{For their honeymoon they went to the seaside.}
\end{quote}

Built on the same picture as English \emph{honeymoon}.

\begin{cousinweb}
\emph{la lune de miel} translates English \emph{honeymoon} word for word: \emph{lune} is inherited from Latin \emph{lūna} \textquotedblleft{}moon\textquotedblright{} and \emph{miel} from Latin \emph{mel} \textquotedblleft{}honey\textquotedblright{} (Lewis \& Short).
\end{cousinweb}

\begin{grammarlens}[title={What you've built}]
One more word for telling, explaining and describing at B1.
\end{grammarlens}

\subsection*{Your turn}
\begin{itemize}
\raggedright
  \item \emph{Say it:} \emph{la lune de miel}
  \item \emph{Say it:} \emph{la lune de miel}, once more
  \item \emph{Say it:} a sentence of your own with \emph{la lune de miel}
  \item \emph{Recall:} say the French for wedding ring, then the French for wedding, then say \emph{la lune de miel} again
\end{itemize}

\begin{tcolorbox}[breakable,colback=teal!4,colframe=teal!35!black,title={Before you move on}]
What does la lune de miel mean? (\textquotedblleft{}Honeymoon\textquotedblright{}.) Say it once more.
\end{tcolorbox}
\section[la cérémonie]{la cérémonie — ceremony}

\label{lesson:FR-C489-ceremonie}



\begin{quote}
Before the new one: say the French for wedding, then the French for honeymoon.
\end{quote}

\subsection*{You'll want to know: la cérémonie}
\textbf{la cérémonie} — \textquotedblleft{}ceremony\textquotedblright{}.

\begin{quote}
La cérémonie a eu lieu dans un petit jardin.
\end{quote}

\begin{quote}
\emph{The ceremony took place in a small garden.}
\end{quote}

\emph{Sans cérémonie} — informally.

\begin{cousinweb}
\emph{la cérémonie} is a learned borrowing from Latin \emph{caerimōnia} \textquotedblleft{}a sacred rite, reverence\textquotedblright{} (EtymDB; Lewis \& Short).
\end{cousinweb}

\begin{grammarlens}[title={What you've built}]
One more word for telling, explaining and describing at B1.
\end{grammarlens}

\subsection*{Your turn}
\begin{itemize}
\raggedright
  \item \emph{Say it:} \emph{la cérémonie}
  \item \emph{Say it:} \emph{la cérémonie}, once more
  \item \emph{Say it:} a sentence of your own with \emph{la cérémonie}
  \item \emph{Recall:} say the French for wedding, then the French for honeymoon, then say \emph{la cérémonie} again
\end{itemize}

\begin{tcolorbox}[breakable,colback=teal!4,colframe=teal!35!black,title={Before you move on}]
What does la cérémonie mean? (\textquotedblleft{}Ceremony\textquotedblright{}.) Say it once more.
\end{tcolorbox}
\section[(dialogue)]{Retour sur les chapitres 482–489}

\label{lesson:FR-R489-revision-482-489}



\begin{quote}
Say the French for to sneeze and allergy.
\end{quote}

\subsection*{Your turn}
Say these aloud:

\begin{itemize}
\raggedright
  \item to sneeze, allergy, pollen, itch, to scratch
  \item nausea, to vomit, migraine, dizziness, to digest
  \item ambulance, first-aider, stretcher, injection, waiting room
  \item to heal, recovery, to restore, convalescence, to look after
  \item insomnia, to snore, nightmare, to yawn, nap
  \item in love, love at first sight, to chat up, tenderness, to stroke
  \item cuddle, to cuddle, closeness, to whisper, to entrust
  \item engagement, wedding ring, wedding, honeymoon, ceremony
\end{itemize}

\begin{tcolorbox}[breakable,colback=teal!4,colframe=teal!35!black,title={Before you move on}]
To sneeze? (\textbf{éternuer}.) Allergy? (\textbf{l'allergie}.) Pollen? (\textbf{le pollen}.) Itch? (\textbf{la démangeaison}.) To scratch? (\textbf{gratter}.) Nausea? (\textbf{la nausée}.) To vomit? (\textbf{vomir}.) Migraine? (\textbf{la migraine}.) Dizziness? (\textbf{le vertige}.) To digest? (\textbf{digérer}.) Ambulance? (\textbf{l'ambulance}.) First-aider? (\textbf{le secouriste}.) Stretcher? (\textbf{le brancard}.) Injection? (\textbf{la piqûre}.) Waiting room? (\textbf{la salle d'attente}.) To heal? (\textbf{guérir}.) Recovery? (\textbf{la guérison}.) To restore? (\textbf{rétablir}.) Convalescence? (\textbf{la convalescence}.) To look after? (\textbf{soigner}.) Insomnia? (\textbf{l'insomnie}.) To snore? (\textbf{ronfler}.) Nightmare? (\textbf{le cauchemar}.) To yawn? (\textbf{bâiller}.) Nap? (\textbf{la sieste}.) In love? (\textbf{amoureux}.) Love at first sight? (\textbf{le coup de foudre}.) To chat up? (\textbf{draguer}.) Tenderness? (\textbf{la tendresse}.) To stroke? (\textbf{caresser}.) Cuddle? (\textbf{le câlin}.) To cuddle? (\textbf{câliner}.) Closeness? (\textbf{la complicité}.) To whisper? (\textbf{chuchoter}.) To entrust? (\textbf{confier}.) Engagement? (\textbf{les fiançailles}.) Wedding ring? (\textbf{l'alliance}.) Wedding? (\textbf{les noces}.) Honeymoon? (\textbf{la lune de miel}.) Ceremony? (\textbf{la cérémonie}.)
\end{tcolorbox}
